\ejercicio{}

Supóngase que no hay ciclos negativos alcanzables desde $s$. Entonces todo camino más corto se puede tomar simple, pues quitar un ciclo de peso no negativo no lo encarece, y un camino simple usa a lo sumo $n - 1$ arcos. Por inducción sobre $i$, después de la ronda $i$, $d[v]$ es a lo sumo el costo del camino más corto de $s$ a $v$ con a lo sumo $i$ arcos: si ese camino termina en el arco $(u, v)$, su prefijo hasta $u$ usa a lo sumo $i - 1$ arcos, y relajar $(u, v)$ en la ronda $i$ deja $d[v]$ por debajo de su costo. Con $i = n - 1$, cada $d[v]$ ya es el costo mínimo al terminar la ronda $n - 1$, y ningún $d[v]$ baja del costo mínimo porque siempre es el costo de algún camino. Así, ninguna relajación de la ronda $n$ disminuye una distancia. Por contrarrecíproca, si alguna disminuye, hay un ciclo negativo alcanzable desde $s$. \hfill$\blacksquare$
